\section{The Undying Worm, Exterior Darkness, and the Fixed Will}

In the Book of Job, the suffering patriarch cries out for a moment's respite before descending ``to a land that is dark and covered with the mist of death: A land of misery and darkness, where the shadow of death, and no order, but everlasting horror dwelleth'' (Job 10:21--22). Reading those words in the ninth book of his \textit{Moralia on Job}, St.\ Gregory the Great pauses over the phrase ``no order'' (\textit{nullus ordo}). God Himself is the author of all order; how can anything in His creation be without order? Gregory answers:

\begin{quote}
Therefore there is done upon the wretches death without death, an end without ending, failing without failing; in that both death lives, and the end is ever beginning, and the failing is unable to fail. Therefore whereas death at the same time slays and does not extinguish, pain torments but does not banish fear, the flame burns but does not dispel the darkness, for all that is gathered from a knowledge of the present life, the punishments are without order, in that they do not retain their own character through all particulars.\endnote{St.\ Gregory the Great, \textit{Morals on the Book of Job} IX.100 (IX.66 in PL 75), Library of the Fathers translation.}
\end{quote}

In this life every affliction obeys a natural limit: death ends pain, an ending closes a process, and a fire that burns also illuminates. In hell those natural bounds are suspended by divine justice, and Scripture places three companions beside the unquenchable fire: the worm that dies not, the exterior darkness, and the fixed anguish of a will hardened against God.

The prophet Isaias closes his entire book with the scene of the new heavens and the new earth and, outside, the transgressors of whom ``their worm shall not die, and their fire shall not be quenched'' (Isa 66:24). The Book of Judith declares that the Lord Almighty ``will give fire, and worms into their flesh, that they may burn, and may feel for ever'' (Jdt 16:21), and Ecclesiasticus counsels humility because ``the vengeance on the flesh of the ungodly is fire and worms'' (Ecclus 7:19). Our Lord Himself takes up Isaias's words three times in Mark 9:43, 45, and 47.

As we have seen, St.\ Augustine in \textit{The City of God} XXI.9 is completely firm about the bodily reality of hell's fire, yet explicitly leaves the reader free on the nature of the worm:

\begin{quote}
Let each man, then, choose which he will, whether he thinks that the worm is real and pertains to the body, or that spiritual things are meant by bodily representations, and that it belongs to the soul. But which of these is true will be more readily discovered by the facts themselves, when there shall be in the saints such knowledge as shall not require that their own experience teach them the nature of these punishments\dots\endnote{St.\ Augustine, \textit{The City of God} XXI.9, trans.\ Dods.}
\end{quote}

Augustine's liberty concerns the worm, not the fire. The Thomistic \textit{Supplement} (q.\ 97, a.\ 2) takes up that liberty and argues cleanly for the spiritual reading. After the general renewal of creation at the Last Judgment, no corruptible animal remains in existence, since only human bodies are ordered to incorruption:

\begin{quote}
Consequently the worm ascribed to the damned must be understood to be not of a corporeal but of a spiritual nature: and this is the remorse of conscience, which is called a worm because it originates from the corruption of sin, and torments the soul, as a corporeal worm born of corruption torments by gnawing.\endnote{\textit{Summa Theologiae}, Supplement, q.\ 97, a.\ 2, trans.\ Fathers of the English Dominican Province.}
\end{quote}

What is it like to be gnawed forever by one's own conscience? St.\ Bernard of Clairvaux, writing to Pope Eugenius III in the fifth book of \textit{De consideratione}, observes that in the final exposure of sin no gaze in heaven or on earth is harder for a darkened conscience to bear or less possible to flee than its own.\endnote{St.\ Bernard of Clairvaux, \textit{De consideratione} V.12.}

The fifth chapter of the Book of Wisdom lets us hear that conscience speaking from hell (\textit{Talia dixerunt in inferno hi qui peccaverunt}, Wis 5:14). Seeing the just stand in constancy among the children of God, the wicked groan in anguish of spirit: ``We fools esteemed their life madness, and their end without honour. \dots\ What hath pride profited us? or what advantage hath the boasting of riches brought us? All those things are passed away like a shadow, and like a post that runneth on, And as a ship that passeth through the waves: whereof when it is gone by, the trace cannot be found'' (Wis 5:4, 8--10). St.\ Alphonsus Liguori, meditating on Mark 9:47 in his twenty-eighth Consideration of \textit{Preparation for Death}, gathers this undying remorse under three bitter heads: the trifles for which the sinner sold eternity, like Esau howling after trading his birthright for lentils (Gen 27:34) or Jonathan tasting a little honey (1 Kings 14:43); the little that was asked in order to be saved, while graces, sacraments, and years of mercy stood open; and the boundless good of God Himself forfeited not by fate or another's malice, but by one's own free hand.\endnote{St.\ Alphonsus Maria de Liguori, \textit{Preparation for Death}, Consideration XXVIII, First--Third Points (Dublin: James Duffy, 1843), pp.\ 340--347.}

Beside the undying worm stands ``the exterior darkness'' (Matt 8:12; 22:13; 25:30)---what St.\ Jude calls ``the storm of darkness'' reserved forever (Jude 1:13). How can a place be filled with fire and yet be dark? St.\ Gregory the Great explains in the \textit{Moralia} that hell's flame gives ``no light for comfort,'' yet keeps just enough gloomy visibility ``that it may torment the more,'' so that the lost behold those companions for whose carnal affection they transgressed against God, just as the rich man in Luke 16 saw Lazarus afar off in Abraham's bosom to the sharp increase of his grief.\endnote{St.\ Gregory the Great, \textit{Morals on the Book of Job} IX.101 (IX.66 in PL 75).} The \textit{Supplement} (q.\ 97, a.\ 4) cites both Gregory and St.\ Basil's comment on Psalm 28:7 (``The voice of the Lord divideth the flame of fire'')---that God divides the brightness of fire for the joy of the blessed from its burning heat for the torment of the wicked---and concludes that ``simply speaking, the place is dark. Yet by Divine disposition, there is a certain amount of light, as much as suffices for seeing those things which are capable of tormenting the soul.''\endnote{\textit{Summa Theologiae}, Supplement, q.\ 97, a.\ 4.} Writing without scholastic terminology about her own mystical vision of hell, St.\ Teresa of Avila reports the same paradox: ``There was no light, but all was thick darkness. I do not understand how it is; though there was no light, yet everything that can give pain by being seen was visible.''\endnote{St.\ Teresa of Avila, \textit{The Life of St.\ Teresa of Jesus} ch.\ 32, \S 5, trans.\ Lewis, rev.\ Zimmerman, p.\ 287.}

Deepest of all, beneath fire, worm, and darkness, lies the obstinacy of the damned will. Why do the damned not simply cry out in true contrition after a thousand years, beg God's pardon, and walk free? The question touches the core of human freedom and death. St.\ John of Damascus formulated the patristic principle quoted throughout the Latin Middle Ages: ``death is to man what their fall was to the angels.''\endnote{St.\ John of Damascus, \textit{Exposition of the Orthodox Faith} II.4, cited in \textit{Summa Theologiae}, Supplement, q.\ 98, a.\ 2, obj.\ 3.} While the soul is united to the mortal body in this passing world, the will remains mobile: grace knocks at every hour, habits can be broken, and a lifelong sinner can turn in a single instant like the thief on the cross. At death, the pilgrimage of testing ends, and the soul's final choice hardens into permanence.

The \textit{Supplement} (q.\ 98, aa.\ 1--2) analyzes that fixed state with precision. Every damned soul retains its \textit{natural} will, which comes from God the Author of nature and still recoils from suffering and longs for happiness in the abstract; that is why the damned never cease to feel the ache of what they have lost. Yet their \textit{deliberate} will---the chosen orientation of the self---is wholly and permanently turned away from the ultimate end. Without direction to God, no deliberate act of the will can be good. Do the damned then repent, as Wisdom 5:3 says (``repenting, and groaning for anguish of spirit'')? Article 2 distinguishes direct from indirect repentance:

\begin{quote}
He repents of a sin directly who hates sin as such: and he repents indirectly who hates it on account of something connected with it, for instance punishment or something of that kind. Accordingly the wicked will not repent of their sins directly, because consent in the malice of sin will remain in them; but they will repent indirectly, inasmuch as they will suffer from the punishment inflicted on them for sin.\endnote{\textit{Summa Theologiae}, Supplement, q.\ 98, a.\ 2.}
\end{quote}

The damned hate the flames, hate the humiliation, and hate the loss, yet they do not hate the pride and malice that placed them there. St.\ Alphonsus draws out the dread consequence: ``Perhaps the miserable man can love God in Hell, and can resign himself to his will? No: if he could do this, Hell should not be Hell. \dots\ his Hell shall consist in the conviction that God is an infinite good, and that he is compelled to hate him, while he sees that he is worthy of infinite love.''\endnote{St.\ Alphonsus Maria de Liguori, \textit{Preparation for Death}, Consideration XXVI, Third Point, p.\ 325.} That is why hell is locked from within as well as from without: it is not a prison holding repentant lovers who long to embrace their Father, but the eternal dwelling of wills that have refused His mercy to the end and remain confirmed in that refusal.
